\documentclass[10pt,a4paper]{amsart}
\usepackage[margin=19mm]{geometry}
\usepackage{amsmath,amssymb,mathtools}
\usepackage[scr=rsfs,cal=euler]{mathalfa}
\usepackage{xfrac}
\usepackage[unicode,hidelinks,bookmarks=false]{hyperref}
\newcommand{\ie}{i.e.\ }
\newcommand{\cf}{cf.\ }
\newcommand{\resp}{respectively\ }
\newcommand{\Subsection}{\subsection}
\providecommand{\nobreakdash}{\nobreak\hbox{-}}
\setlength{\emergencystretch}{2em}
\multlinegap0pt
\usepackage{amssymb}
\newtheorem{lem}{Lemma}
\DeclareMathOperator{\hdim}{hdim}
\def\k {\mathrm k}
\title{Critical subgraphs and the regularity of symbolic powers of cover ideals of graphs}
\author{Nguyen Thu Hang, Thanh Vu}
\date{Source: arXiv:2605.20603v1 (CC BY 4.0)}
\begin{document}
\maketitle

\noindent\emph{Source and licence.} Critical subgraphs and the regularity of symbolic powers of cover ideals of graphs. Version arXiv:2605.20603v1 (CC BY 4.0), distributed by arXiv under the Creative Commons Attribution 4.0 International licence, http://creativecommons.org/licenses/by/4.0/, as recorded in the arXiv licence metadata for this exact version and shown on the abstract page https://arxiv.org/abs/2605.20603v1. Full licence notice: \texttt{licenses/arxiv-2605.20603-CC-BY-4.0.txt}.\par\smallskip

\noindent\textit{Editorial scope.} Counted unit: the article's lem lem\_reduction\_1 (source0.tex lines 782--794). The article's complete original proof of it is delivered with it (source0.tex lines 796--876, one \texttt{proof} environment of 81 lines). No mathematics is written, repaired or re-derived: the statement and the proof are the article's own lines, in the article's own notation, and no retained line is substituted or re-flowed.  No further source-local context is needed: every cross-reference of every retained line resolves inside the two delivered spans.  Nothing was omitted from any retained span, so the package carries no omission record.  No retained line contains a citation, so the wrapper carries no bibliography and no bibliography entry.  No retained line contains a figure environment, an image-inclusion command or drawing code, so no image is delivered and the package declares no asset dependency.  Editorial formatting only, in addition: an AMS \texttt{amsart} preamble that restates the article's own declarations and macros used by the retained lines, namely the environments \texttt{lem} on separate counters, together with the standard packages those lines need.  The retained environments print in this document as Lemma 1 in order of appearance, whereas the article prints the same items with the numbering of its own full document.  The complete native source of the article as distributed in the arXiv e-print for this exact version is bundled unchanged as \texttt{sources/200999/source0.tex}.
\begin{lem}\label{lem_reduction_1}
Let $G$ be a simple graph, and let $u$ be a vertex of $G$. Assume that $u$ has at least two leaf neighbors $v_1$ and $v_2$. Let $H = G \setminus v_1$. Then:
\begin{enumerate}
    \item
    \(    \tau_{\max}(H) \le \tau_{\max}(G) - 1;    \)
    
    \item
    \(    \hdim_{\k}\bigl(\Delta(J(H))\bigr)
    =
    \hdim_{\k}\bigl(\Delta(J(G))\bigr) - 1.
    \)
\end{enumerate}
\end{lem}

\begin{proof}
\textbf{(1)} Let $C$ be a minimal cover of $H$ with $|C| = \tau_{\max}(H)$. If $u \notin C$, then $C \cup \{v_1\}$ is a minimal cover of $G$. Hence,
\[
\tau_{\max}(G) \ge |C| + 1 = \tau_{\max}(H) + 1,
\]
which proves the claim. Now assume that $u \in C$. Let
\[
L = \{v \in N_H(u) \mid v \text{ is a leaf}\}
 \text{ and } 
A = \{v \in N_H(u) \mid v \notin C \text{ and } v \text{ is not a leaf}\}.
\]
Then
\[
C' = (C \setminus \{u\}) \cup L \cup A
\]
is a minimal cover of $H$, and \(|C'| \ge |C|.\) Since $u \notin C'$, we are reduced to the previous case. Therefore,
\[
\tau_{\max}(H) \le \tau_{\max}(G) - 1.
\]

\medskip

\textbf{(2)} Let
\[
\Delta = \Delta(J(G)),
\qquad
F = V(G) \setminus \{u,v_1\},
\qquad
\Gamma = \Delta \setminus F.
\]
Then $\Gamma$ is a cone with apex $v_1$, and hence is acyclic. By the Mayer--Vietoris sequence, we obtain
\begin{equation}\label{eq1}
\hdim_{\k}(\Delta)
=
\hdim_{\k}(F \cap \Gamma) + 1.
\end{equation}
Now, $F \cap \Gamma$ is the union of the facet
\[
F_1 = V(G) \setminus \{u,v_1,v_2\},
\]
and a simplicial complex $\Gamma_1$ whose facets are of the form
\[
V(G) \setminus \bigl(\{u,v_1\} \cup e\bigr),
\]
where $e$ is an edge of $H$ other than $\{u,v_2\}$. Note that $\Gamma_1$ is a cone with apex $v_2$. Hence,
\begin{equation}\label{eq2}
\hdim_\k (F \cap \Gamma) = \hdim_{\k}(F_1 \cup \Gamma_1)
=
\hdim_{\k}(F_1 \cap \Gamma_1) + 1.
\end{equation}
Now, $F_1 \cap \Gamma_1$ is the simplicial complex $\Sigma$ consisting of the facets
\[
V(G) \setminus \bigl(\{u,v_1,v_2\} \cup e\bigr),
\]
where $e$ is an edge of $H$ other than $\{u,v_2\}$.

On the other hand, decompose $\Delta(J(H))$ as \( \Delta(J(H)) = F_2 \cup \Lambda,\) where
\[
F_2 = V(H) \setminus \{u,v_2\}
     = V(G) \setminus \{u,v_1,v_2\},
\]
and $\Lambda$ is the subcomplex generated by all facets other than $F_2$. Then
\[
F_2 \cap \Lambda = \Sigma.
\]
Since $\Lambda$ is also a cone with apex $v_2$, it follows that
\[
\hdim_{\k}\bigl(\Delta(J(H))\bigr)
=
\hdim_{\k}(F_2 \cap \Lambda) + 1
=
\hdim_{\k}(\Sigma) + 1.
\]
Combining this with Eq.~\eqref{eq1} and Eq.~\eqref{eq2}, we conclude that
\[
\hdim_{\k}\bigl(\Delta(J(H))\bigr)
=
\hdim_{\k}\bigl(\Delta(J(G))\bigr) - 1.
\]
This completes the proof.
\end{proof}

\end{document}
